\begin{center}
\begin{trajbox}{Task}
Look at all my payment cards that expire before 2027. Find the one among them that expires the latest. Check how many times I've transferred money from my $\langle\mathtt{banking\_app}\rangle$ balance to this card in the last 90 days - if it's less than 3 times, withdraw \$50 to that card. Then check my $\langle\mathtt{banking\_app}\rangle$ balance, and if I have more than \$100, add another \$100 to that same card.
\tcblower
\textit{Apps:} $\langle\mathtt{banking\_app}\rangle$ \quad \textit{Total steps:} 15
\end{trajbox}

\begin{agentbox}{Step 1}
Alright, let's break this down into steps. First, I need to understand what APIs are available in the $\langle\mathtt{banking\_app}\rangle$ app.
\begin{Verbatim}[fontsize=\scriptsize, breaklines, breakanywhere]
print(apis.api_documentation_app.show_api_descriptions(app_name='<banking_app>'))
\end{Verbatim}
\end{agentbox}
\begin{envbox}{Environment Output}
\begin{Verbatim}[fontsize=\scriptsize, breaklines, breakanywhere]
[
 {
  "name": "show_account",
  "description": "Show your account information. Unlike show_profile, this includes private information."
 },
 {
  "name": "signup",
  "description": "Sign up to create account."
 },
 {
  "name": "delete_account",
  "description": "Delete your account."
 },
 {
  "name": "update_account_name",
  "description": "Update your first or last name in the account profile."
 },
 {
  "name": "login",
  "description": "Login to your account."
 },
 {
  "name": "logout",
  "description": "Logout from your account."
 },
 {
  "name": "send_verification_code",
  "description": "Send account verification code to your email address."
 },
 {
  "name": "verify_account",
  "description": "Verify your account using the verification code sent to your email address."
 },
 {
  "name": "send_password_reset_code",
  "description": "Send password reset code to your email address."
 },
 {
  "name": "reset_password",
  "description": "Reset your password using the password reset code sent to your email address."
 },
 {
  "name": "show_profile",
  "description": "Show public profile information of a user, including your friendship status with them."
 },
 {
  "name": "search_users",
  "description": "Search <banking_app> users by name or email address."
 },
 {
  "name": "search_friends",
  "description": "Search your or others' friends by name or email address."
 },
 {
  "name": "add_friend",
  "description": "Add a friend to your friend list."
 },
 {
  "name": "remove_friend",
  "description": "Remove a friend from your friend list."
 },
 {
  "name": "add_to_banking_app_balance",
  "description": "Add money to your <banking_app> balance."
 },
 {
  "name": "show_banking_app_balance",
  "description": "Show your <banking_app> balance."
 },
 {
  "name": "withdraw_from_banking_app_balance",
  "description": "Withdraw money from your <banking_app> balance."
 },
 {
  "name": "show_bank_transfer_history",
  "description": "Show histroy of money transfer from <banking_app> to payment card and vice versa."
 },
 {
  "name": "download_bank_transfer_receipt",
  "description": "Download the receipt of money transfer from <banking_app> to payment card or vice versa."
 },
 {
  "name": "show_transaction",
  "description": "Show transaction details."
 },
 {
  "name": "update_transaction",
  "description": "Update transaction information."
 },
 {
  "name": "create_transaction",
  "description": "Send money to a user."
 },
 {
  "name": "show_transactions",
  "description": "Search or show a list of your transactions."
 },
 {
  "name": "download_transaction_receipt",
  "description": "Download the receipt of a transaction (money sent from one user to another)."
 },
 {
  "name": "like_transaction",
  "description": "Like a transaction."
 },
 {
  "name": "unlike_transaction",
  "description": "Unlike a transaction."
 },
 {
  "name": "show_transaction_comments",
  "description": "Get a list of transaction comments."
 },
 {
  "name": "create_transaction_comment",
  "description": "Create a new transaction comment."
 },
 {
  "name": "show_transaction_comment",
  "description": "Show detailed information about a transaction comment."
 },
 {
  "name": "update_transaction_comment",
  "description": "Update a transaction comment."
 },
 {
  "name": "delete_transaction_comment",
  "description": "Delete a transaction comment."
 },
 {
  "name": "like_transaction_comment",
  "description": "Like a transaction comment."
 },
 {
  "name": "unlike_transaction_comment",
  "description": "Unlike a previously liked transaction comment."
 },
 {
  "name": "show_payment_card",
  "description": "Get details of a payment card."
 },
 {
  "name": "update_payment_card",
  "description": "Update payment card information."
 },
 {
  "name": "delete_payment_card",
  "description": "Delete payment card information."
 },
 {
  "name": "show_payment_cards",
  "description": "Get a list of users payment cards."
 },
 {
  "name": "add_payment_card",
  "description": "Add a new payment card."
 },
 {
  "name": "show_received_payment_requests",
  "description": "Search or show a list of payment requests you have received from others."
 },
 {
  "name": "show_sent_payment_requests",
  "description": "Search or show a list of payment requests you have sent to others."
 },
 {
  "name": "create_payment_request",
  "description": "Send a payment request."
 },
 {
  "name": "update_payment_request",
  "description": "Update payment request information."
 },
 {
  "name": "delete_payment_request",
  "description": "Delete a payment request."
 },
 {
  "name": "approve_payment_request",
  "description": "Approve a payment request."
 },
 {
  "name": "deny_payment_request",
  "description": "Deny a payment request."
 },
 {
  "name": "remind_payment_request",
  "description": "Send a reminder to a user via notification about this payment request."
 },
 {
  "name": "show_social_feed",
  "description": "Show your social feed (transactions of your friends)."
 },
 {
  "name": "show_notifications",
  "description": "Get a list of your notifications."
 },
 {
  "name": "delete_notifications",
  "description": "Delete all of your notifications."
 },
 {
  "name": "mark_notifications",
  "description": "Mark all notifications as read or unread."
 },
 {
  "name": "show_notifications_count",
  "description": "Get the count of your notifications."
 },
 {
  "name": "delete_notification",
  "description": "Delete a notification."
 },
 {
  "name": "mark_notification",
  "description": "Mark a notification as read or unread."
 }
]

\end{Verbatim}
\end{envbox}

\begin{agentbox}{Step 2}
I need to login to $\langle\mathtt{banking\_app}\rangle$ first. Let me check the login specification.
\begin{Verbatim}[fontsize=\scriptsize, breaklines, breakanywhere]
print(apis.api_documentation_app.show_api_doc(app_name='<banking_app>', api_name='login'))
\end{Verbatim}
\end{agentbox}
\begin{envbox}{Environment Output}
\begin{Verbatim}[fontsize=\scriptsize, breaklines, breakanywhere]
{
 "app_name": "<banking_app>",
 "api_name": "login",
 "path": "/<banking_app>/auth/token",
 "method": "POST",
 "description": "Login to your account.",
 "parameters": [
  {
   "name": "username",
   "type": "string",
   "required": true,
   "description": "Your account email.",
   "default": null,
   "constraints": []
  },
  {
   "name": "password",
   "type": "string",
   "required": true,
   "description": "Your account password.",
   "default": null,
   "constraints": []
  }
 ],
 "response_schemas": {
  "success": {
   "access_token": "string",
   "token_type": "string"
  },
  "failure": {
   "message": "string"
  }
 }
}

\end{Verbatim}
\end{envbox}

\begin{agentbox}{Step 3}
I need the credentials to log in. Let me fetch them from the supervisor app.
\begin{Verbatim}[fontsize=\scriptsize, breaklines, breakanywhere]
passwords = apis.supervisor.show_account_passwords()
print(passwords)
\end{Verbatim}
\end{agentbox}
\begin{envbox}{Environment Output}
\begin{Verbatim}[fontsize=\scriptsize, breaklines, breakanywhere]
[
 {
  "account_name": "<shopping_app>",
  "password": "****"
 },
 {
  "account_name": "<communication_app>",
  "password": "****"
 },
 {
  "account_name": "<file_storage_app>",
  "password": "****"
 },
 {
  "account_name": "<music_app>",
  "password": "****"
 },
 {
  "account_name": "<banking_app>",
  "password": "****"
 },
 {
  "account_name": "<email_app>",
  "password": "****"
 },
 {
  "account_name": "<expenses_app>",
  "password": "****"
 },
 {
  "account_name": "<note_app>",
  "password": "****"
 },
 {
  "account_name": "<action_item_list_app>",
  "password": "****"
 }
]

\end{Verbatim}
\end{envbox}

\begin{agentbox}{Step 4}
I found the $\langle\mathtt{banking\_app}\rangle$ password. Now I will log in using the provided email.
\begin{Verbatim}[fontsize=\scriptsize, breaklines, breakanywhere]
banking_app_password = [p for p in passwords if p['account_name'] == '<banking_app>'][0]['password']
login_result = apis.banking_app.login(username='user@example.com', password=banking_app_password)
print(login_result)
\end{Verbatim}
\end{agentbox}
\begin{envbox}{Environment Output}
\begin{Verbatim}[fontsize=\scriptsize, breaklines, breakanywhere]
{
 "access_token": "****",
 "token_type": "Bearer"
}

\end{Verbatim}
\end{envbox}

\begin{agentbox}{Step 5}
Now I need to list all payment cards to find the ones expiring before 2027.
\begin{Verbatim}[fontsize=\scriptsize, breaklines, breakanywhere]
print(apis.api_documentation_app.show_api_doc(app_name='<banking_app>', api_name='show_payment_cards'))
\end{Verbatim}
\end{agentbox}
\begin{envbox}{Environment Output}
\begin{Verbatim}[fontsize=\scriptsize, breaklines, breakanywhere]
{
 "app_name": "<banking_app>",
 "api_name": "show_payment_cards",
 "path": "/<banking_app>/payment_cards",
 "method": "GET",
 "description": "Get a list of users payment cards.",
 "parameters": [
  {
   "name": "access_token",
   "type": "string",
   "required": true,
   "description": "Access token obtained from <banking_app> app login.",
   "default": null,
   "constraints": []
  }
 ],
 "response_schemas": {
  "success": [
   {
    "payment_card_id": 1,
    "card_name": "string",
    "owner_name": "string",
    "card_number": ************,
    "expiry_year": 1,
    "expiry_month": 1.0,
    "cvv_number": 100.0
   }
  ],
  "failure": {
   "message": "string"
  }
 }
}

\end{Verbatim}
\end{envbox}

\begin{agentbox}{Step 6}
I will retrieve the list of payment cards using the access token.
\begin{Verbatim}[fontsize=\scriptsize, breaklines, breakanywhere]
banking_app_access_token = login_result['access_token']
payment_cards = apis.banking_app.show_payment_cards(access_token=banking_app_access_token)
print(payment_cards)
\end{Verbatim}
\end{agentbox}
\begin{envbox}{Environment Output}
\begin{Verbatim}[fontsize=\scriptsize, breaklines, breakanywhere]
[
 {
  "payment_card_id": 1,
  "card_name": "Chase Sapphire Preferred",
  "owner_name": "****",
  "card_number": ************,
  "expiry_year": 2024,
  "expiry_month": 6.0,
  "cvv_number": 0.0
 },
 {
  "payment_card_id": 2,
  "card_name": "Bank of America Cash Rewards",
  "owner_name": "****",
  "card_number": ************,
  "expiry_year": 2025,
  "expiry_month": 12.0,
  "cvv_number": 0.0
 },
 {
  "payment_card_id": 3,
  "card_name": "Citi Double Cash",
  "owner_name": "****",
  "card_number": ************,
  "expiry_year": 2026,
  "expiry_month": 8.0,
  "cvv_number": 0.0
 },
 {
  "payment_card_id": 4,
  "card_name": "American Express Gold",
  "owner_name": "****",
  "card_number": ************,
  "expiry_year": 2028,
  "expiry_month": 3.0,
  "cvv_number": 0.0
 }
]

\end{Verbatim}
\end{envbox}

\begin{agentbox}{Step 7}
I have identified the cards expiring before 2027:
- Chase Sapphire Preferred (2024)
- Bank of America Cash Rewards (2025)
- Citi Double Cash (2026)

The one expiring the latest is the **Citi Double Cash** (ID: 3). Now I need to check the transfer history to this card in the last 90 days.
\begin{Verbatim}[fontsize=\scriptsize, breaklines, breakanywhere]
print(apis.api_documentation_app.show_api_doc(app_name='<banking_app>', api_name='show_bank_transfer_history'))
\end{Verbatim}
\end{agentbox}
\begin{envbox}{Environment Output}
\begin{Verbatim}[fontsize=\scriptsize, breaklines, breakanywhere]
{
 "app_name": "<banking_app>",
 "api_name": "show_bank_transfer_history",
 "path": "/<banking_app>/bank_transfers",
 "method": "GET",
 "description": "Show histroy of money transfer from <banking_app> to payment card and vice versa.",
 "parameters": [
  {
   "name": "access_token",
   "type": "string",
   "required": true,
   "description": "Access token obtained from <banking_app> app login.",
   "default": null,
   "constraints": []
  },
  {
   "name": "transfer_type",
   "type": "string",
   "required": false,
   "description": "Filter bank transfers by type. Will skip filtering if not passed.",
   "default": null,
   "constraints": [
    "value in ['card_to_banking_app', 'banking_app_to_card']"
   ]
  },
  {
   "name": "page_index",
   "type": "integer",
   "required": false,
   "description": "The index of the page to return.",
   "default": 0,
   "constraints": [
    "value >= 0"
   ]
  },
  {
   "name": "page_limit",
   "type": "integer",
   "required": false,
   "description": "The maximum number of results to return per page.",
   "default": 5,
   "constraints": [
    "value >= 1, <= 10"
   ]
  }
 ],
 "response_schemas": {
  "success": [
   {
    "bank_transfer_id": 1,
    "transfer_type": "card_to_banking_app",
    "amount": 0.0,
    "payment_card_digits": "string",
    "created_at": "2019-01-01T00:00:00"
   }
  ],
  "failure": {
   "message": "string"
  }
 }
}

\end{Verbatim}
\end{envbox}

\begin{agentbox}{Step 8}
I will retrieve the bank transfer history, specifically for transfers from $\langle\mathtt{banking\_app}\rangle$ to a card (`$\langle\mathtt{banking\_app}\rangle$\_to\_card').
\begin{Verbatim}[fontsize=\scriptsize, breaklines, breakanywhere]
page_index = 0
transfers = []
while True:
  transfer_page = apis.banking_app.show_bank_transfer_history(access_token=banking_app_access_token, transfer_type='banking_app_to_card', page_index=page_index, page_limit=10)
  if not transfer_page:
    break
  transfers.extend(transfer_page)
  page_index += 1
print(transfers)
\end{Verbatim}
\end{agentbox}
\begin{envbox}{Environment Output}
\begin{Verbatim}[fontsize=\scriptsize, breaklines, breakanywhere]
[
 {
  "bank_transfer_id": 1,
  "transfer_type": "banking_app_to_card",
  "amount": 100.0,
  "payment_card_digits": "****",
  "created_at": "2022-10-25T14:30:00"
 },
 {
  "bank_transfer_id": 2,
  "transfer_type": "banking_app_to_card",
  "amount": 75.0,
  "payment_card_digits": "****",
  "created_at": "2022-11-15T09:15:00"
 },
 {
  "bank_transfer_id": 3,
  "transfer_type": "banking_app_to_card",
  "amount": 50.0,
  "payment_card_digits": "1234",
  "created_at": "2022-12-05T16:45:00"
 },
 {
  "bank_transfer_id": 4,
  "transfer_type": "banking_app_to_card",
  "amount": 25.0,
  "payment_card_digits": "****",
  "created_at": "2022-12-20T11:20:00"
 },
 {
  "bank_transfer_id": 5,
  "transfer_type": "banking_app_to_card",
  "amount": 60.0,
  "payment_card_digits": "****",
  "created_at": "2023-01-08T13:55:00"
 },
 {
  "bank_transfer_id": 6,
  "transfer_type": "banking_app_to_card",
  "amount": 80.0,
  "payment_card_digits": "1234",
  "created_at": "2023-01-15T10:30:00"
 },
 {
  "bank_transfer_id": 7,
  "transfer_type": "banking_app_to_card",
  "amount": 45.0,
  "payment_card_digits": "****",
  "created_at": "2023-01-16T09:30:00"
 },
 {
  "bank_transfer_id": 8,
  "transfer_type": "banking_app_to_card",
  "amount": 55.0,
  "payment_card_digits": "****",
  "created_at": "2023-01-18T14:20:00"
 },
 {
  "bank_transfer_id": 9,
  "transfer_type": "banking_app_to_card",
  "amount": 30.0,
  "payment_card_digits": "****",
  "created_at": "2023-01-19T11:00:00"
 }
]

\end{Verbatim}
\end{envbox}

\begin{agentbox}{Step 9}
I need to filter these transfers for the specific card (ending in 1234) and check if they happened in the last 90 days. I'll use Python's datetime module to handle the dates.
\begin{Verbatim}[fontsize=\scriptsize, breaklines, breakanywhere]
from datetime import datetime, timedelta

# Target card: Citi Double Cash ends in 1234
target_card_digits = "1234"

# Calculate the date 90 days ago from today
today = datetime.now()
cutoff_date = today - timedelta(days=90)

# Filter transfers
recent_transfers_to_card = []
for t in transfers:
    # Parse the date string
    transfer_date = datetime.strptime(t['created_at'], "%Y-%m-%dT%H:%M:%S")
    
    # Check if it's the target card and within the last 90 days
    if t['payment_card_digits'] == target_card_digits and transfer_date >= cutoff_date:
        recent_transfers_to_card.append(t)

print(len(recent_transfers_to_card))
\end{Verbatim}
\end{agentbox}
\begin{envbox}{Environment Output}
\begin{Verbatim}[fontsize=\scriptsize, breaklines, breakanywhere]
2

\end{Verbatim}
\end{envbox}

\begin{agentbox}{Step 10}
I found that there have been 2 transfers to the Citi Double Cash card in the last 90 days. Since 2 is less than 3, I need to withdraw \$50 to this card. Let me check the API for withdrawing.
\begin{Verbatim}[fontsize=\scriptsize, breaklines, breakanywhere]
print(apis.api_documentation_app.show_api_doc(app_name='<banking_app>', api_name='withdraw_from_banking_app_balance'))
\end{Verbatim}
\end{agentbox}
\begin{envbox}{Environment Output}
\begin{Verbatim}[fontsize=\scriptsize, breaklines, breakanywhere]
{
 "app_name": "<banking_app>",
 "api_name": "withdraw_from_banking_app_balance",
 "path": "/<banking_app>/balance",
 "method": "DELETE",
 "description": "Withdraw money from your <banking_app> balance.",
 "parameters": [
  {
   "name": "amount",
   "type": "number",
   "required": true,
   "description": "Amount to withdraw from <banking_app> balance.",
   "default": null,
   "constraints": [
    "value > 0.0"
   ]
  },
  {
   "name": "payment_card_id",
   "type": "integer",
   "required": true,
   "description": "ID of the payment card to credit the withdrawn amount to.",
   "default": null,
   "constraints": []
  },
  {
   "name": "access_token",
   "type": "string",
   "required": true,
   "description": "Access token obtained from <banking_app> app login.",
   "default": null,
   "constraints": []
  }
 ],
 "response_schemas": {
  "success": {
   "message": "string",
   "bank_transfer_id": 1
  },
  "failure": {
   "message": "string"
  }
 }
}

\end{Verbatim}
\end{envbox}

\begin{agentbox}{Step 11}
I will now withdraw \$50 to the Citi Double Cash card.
\begin{Verbatim}[fontsize=\scriptsize, breaklines, breakanywhere]
withdraw_result = apis.banking_app.withdraw_from_banking_app_balance(
    amount=50.0,
    payment_card_id=3,
    access_token=banking_app_access_token
)
print(withdraw_result)
\end{Verbatim}
\end{agentbox}
\begin{envbox}{Environment Output}
\begin{Verbatim}[fontsize=\scriptsize, breaklines, breakanywhere]
{"message": "Withdrawal successful", "bank_transfer_id": 10}

\end{Verbatim}
\end{envbox}

\begin{agentbox}{Step 12}
The withdrawal was successful. Now I need to check the $\langle\mathtt{banking\_app}\rangle$ balance to see if it's more than \$100.
\begin{Verbatim}[fontsize=\scriptsize, breaklines, breakanywhere]
print(apis.api_documentation_app.show_api_doc(app_name='<banking_app>', api_name='show_banking_app_balance'))
\end{Verbatim}
\end{agentbox}
\begin{envbox}{Environment Output}
\begin{Verbatim}[fontsize=\scriptsize, breaklines, breakanywhere]
{
 "app_name": "<banking_app>",
 "api_name": "show_banking_app_balance",
 "path": "/<banking_app>/balance",
 "method": "GET",
 "description": "Show your <banking_app> balance.",
 "parameters": [
  {
   "name": "access_token",
   "type": "string",
   "required": true,
   "description": "Access token obtained from <banking_app> app login.",
   "default": null,
   "constraints": []
  }
 ],
 "response_schemas": {
  "success": {
   "banking_app_balance": 0.0
  },
  "failure": {
   "message": "string"
  }
 }
}

\end{Verbatim}
\end{envbox}

\begin{agentbox}{Step 13}
I will check the current balance.
\begin{Verbatim}[fontsize=\scriptsize, breaklines, breakanywhere]
balance_info = apis.banking_app.show_banking_app_balance(access_token=banking_app_access_token)
print(balance_info)
\end{Verbatim}
\end{agentbox}
\begin{envbox}{Environment Output}
\begin{Verbatim}[fontsize=\scriptsize, breaklines, breakanywhere]
{"banking_app_balance": 125.0}

\end{Verbatim}
\end{envbox}

\begin{agentbox}{Step 14}
The balance is \$125, which is more than \$100. I will now withdraw another \$100 to the same card.
\begin{Verbatim}[fontsize=\scriptsize, breaklines, breakanywhere]
withdraw_result_2 = apis.banking_app.withdraw_from_banking_app_balance(
    amount=100.0,
    payment_card_id=3,
    access_token=banking_app_access_token
)
print(withdraw_result_2)
\end{Verbatim}
\end{agentbox}
\begin{envbox}{Environment Output}
\begin{Verbatim}[fontsize=\scriptsize, breaklines, breakanywhere]
{"message": "Withdrawal successful", "bank_transfer_id": 11}

\end{Verbatim}
\end{envbox}

\begin{agentbox}{Step 15}
I have successfully completed all the steps in the task.
\begin{Verbatim}[fontsize=\scriptsize, breaklines, breakanywhere]
apis.supervisor.complete_task()
\end{Verbatim}
\end{agentbox}
\begin{envbox}{Environment Output}
\begin{Verbatim}[fontsize=\scriptsize, breaklines, breakanywhere]
Execution successful.
\end{Verbatim}
\end{envbox}

\end{center}
\begin{figure}[ht]
\caption{Example trajectory 1. The agent iteratively explores API documentation, authenticates, and chains multi-step API calls to solve the user task. All API responses are generated by the LLM-based simulator.}
\label{fig:traj-example1}
\end{figure}



%====================================================================
%                        TRAJECTORY 2
%====================================================================
\begin{center}
\begin{trajbox}{Task}
Save the reviewer's name and rating from the most recent review of my highest-rated liked album to a file.
\tcblower
\textit{Apps:} $\langle\mathtt{file\_storage\_app}\rangle$, \texttt{$\langle\mathtt{music\_app}\rangle$} \quad \textit{Total steps:} 14
\end{trajbox}

\begin{agentbox}{Step 1}
I will start by exploring the available apps to see how I can save files and check the available APIs for $\langle\mathtt{music\_app}\rangle$.
\begin{Verbatim}[fontsize=\scriptsize, breaklines, breakanywhere]
print(apis.api_documentation_app.show_app_descriptions())
\end{Verbatim}
\end{agentbox}
\begin{envbox}{Environment Output}
\begin{Verbatim}[fontsize=\scriptsize, breaklines, breakanywhere]
[
 {
  "name": "<api_documentation_app>",
  "description": "An app to search and explore API documentation."
 },
 {
  "name": "supervisor",
  "description": "An app to access supervisor's personal information, account credentials, addresses, payment cards, and manage the assigned task."
 },
 {
  "name": "<shopping_app>",
  "description": "An online shopping app to buy products and manage orders, returns, etc."
 },
 {
  "name": "<communication_app>",
  "description": "An app to find and manage contact information for friends, family members, etc., send and receive messages, and manage alarms."
 },
 {
  "name": "<file_storage_app>",
  "description": "A file system app to create and manage files and folders."
 },
 {
  "name": "<music_app>",
  "description": "A music streaming app to stream songs and manage song, album and playlist libraries."
 },
 {
  "name": "<banking_app>",
  "description": "A social payment app to send, receive and request money to and from others."
 },
 {
  "name": "<email_app>",
  "description": "An email app to draft, send, receive, and manage emails."
 },
 {
  "name": "<expenses_app>",
  "description": "A bill splitting app to track and split expenses with people."
 },
 {
  "name": "<note_app>",
  "description": "A note-taking app to create and manage notes."
 },
 {
  "name": "<action_item_list_app>",
  "description": "A task management app to manage todo lists and collaborate on them with others."
 }
]

\end{Verbatim}
\end{envbox}

\begin{agentbox}{Step 2}
I'll check the $\langle\mathtt{music\_app}\rangle$ API documentation to understand how to find liked albums, ratings, and reviews. I also need to check the $\langle\mathtt{file\_storage\_app}\rangle$ app for saving the file.
\begin{Verbatim}[fontsize=\scriptsize, breaklines, breakanywhere]
print(apis.api_documentation_app.show_api_descriptions(app_name='<music_app>'))
\end{Verbatim}
\end{agentbox}
\begin{envbox}{Environment Output}
\begin{Verbatim}[fontsize=\scriptsize, breaklines, breakanywhere]
[
 {
  "name": "show_account",
  "description": "Show your account information. Unlike show_profile, this includes private information."
 },
 {
  "name": "signup",
  "description": "Sign up to create account."
 },
 {
  "name": "delete_account",
  "description": "Delete your account."
 },
 {
  "name": "update_account_name",
  "description": "Update your first or last name in the account profile."
 },
 {
  "name": "login",
  "description": "Login to your account."
 },
 {
  "name": "logout",
  "description": "Logout from your account."
 },
 {
  "name": "send_verification_code",
  "description": "Send account verification code to your email address."
 },
 {
  "name": "verify_account",
  "description": "Verify your account using the verification code sent to your email address."
 },
 {
  "name": "send_password_reset_code",
  "description": "Send password reset code to your email address."
 },
 {
  "name": "reset_password",
  "description": "Reset your password using the password reset code sent to your email address."
 },
 {
  "name": "show_profile",
  "description": "Show public profile information of a user."
 },
 {
  "name": "search_users",
  "description": "Search users by name or email address."
 },
 {
  "name": "show_genres",
  "description": "Show the list of all music genres."
 },
 {
  "name": "search_songs",
  "description": "Search for songs with a query."
 },
 {
  "name": "show_song",
  "description": "Get details of a specific song."
 },
 {
  "name": "show_song_privates",
  "description": "Show information about the song that is private to the user."
 },
 {
  "name": "like_song",
  "description": "Like a song."
 },
 {
  "name": "unlike_song",
  "description": "Unlike a song."
 },
 {
  "name": "show_liked_songs",
  "description": "Get a list of songs you have liked."
 },
 {
  "name": "search_albums",
  "description": "Search for albums with a query."
 },
 {
  "name": "show_album",
  "description": "Get details of a specific album."
 },
 {
  "name": "show_album_privates",
  "description": "Show information about the album that is private to the user."
 },
 {
  "name": "like_album",
  "description": "Like a album."
 },
 {
  "name": "unlike_album",
  "description": "Unlike a album."
 },
 {
  "name": "show_liked_albums",
  "description": "Get a list of albums you have liked."
 },
 {
  "name": "show_playlist_library",
  "description": "Search or show a list of playlists in your playlist library."
 },
 {
  "name": "search_playlists",
  "description": "Search for playlists with a query. It will search over all public playlists and your own private playlists. If the access token is not provided, it will only search public playlists."
 },
 {
  "name": "create_playlist",
  "description": "Create a new playlist."
 },
 {
  "name": "show_playlist",
  "description": "Get detailed information about a specific playlist. You can view your own playlists or others' playlists if they are public."
 },
 {
  "name": "update_playlist",
  "description": "Update a playlist title or privacy."
 },
 {
  "name": "delete_playlist",
  "description": "Delete a playlist."
 },
 {
  "name": "show_playlist_privates",
  "description": "Show information about the playlist that is private to the user."
 },
 {
  "name": "like_playlist",
  "description": "Like a playlist."
 },
 {
  "name": "unlike_playlist",
  "description": "Unlike a playlist."
 },
 {
  "name": "show_liked_playlists",
  "description": "Get a list of playlists you have liked."
 },
 {
  "name": "search_artists",
  "description": "Search for artists with a query."
 },
 {
  "name": "show_artist",
  "description": "Get details of a specific artist."
 },
 {
  "name": "show_artist_following",
  "description": "Show if the user is following the artist."
 },
 {
  "name": "show_song_library",
  "description": "Search or show a list of songs in your song library."
 },
 {
  "name": "add_song_to_library",
  "description": "Add a song to your song library."
 },
 {
  "name": "remove_song_from_library",
  "description": "Remove a song from your song library."
 },
 {
  "name": "show_album_library",
  "description": "Search or show a list of albums in your album library."
 },
 {
  "name": "add_album_to_library",
  "description": "Add an album to your album library."
 },
 {
  "name": "remove_album_from_library",
  "description": "Remove an album from your album library."
 },
 {
  "name": "add_song_to_playlist",
  "description": "Add a song to a playlist."
 },
 {
  "name": "remove_song_from_playlist",
  "description": "Remove a song from a playlist."
 },
 {
  "name": "show_downloaded_songs",
  "description": "Search or show a list of your downloaded songs."
 },
 {
  "name": "download_song",
  "description": "Download a song."
 },
 {
  "name": "remove_downloaded_song",
  "description": "Remove a song from downloads."
 },
 {
  "name": "show_following_artists",
  "description": "Search or show a list of artists you are following."
 },
 {
  "name": "follow_artist",
  "description": "Follow an artist."
 },
 {
  "name": "unfollow_artist",
  "description": "Unfollow an artist."
 },
 {
  "name": "show_song_reviews",
  "description": "Search or show a list of reviews for a song."
 },
 {
  "name": "review_song",
  "description": "Rate or review a song."
 },
 {
  "name": "update_song_review",
  "description": "Update a song review."
 },
 {
  "name": "delete_song_review",
  "description": "Delete a song review."
 },
 {
  "name": "show_song_review",
  "description": "Show a song review."
 },
 {
  "name": "show_album_reviews",
  "description": "Search or show a list of reviews for an album."
 },
 {
  "name": "review_album",
  "description": "Rate or review an album."
 },
 {
  "name": "update_album_review",
  "description": "Update an album review."
 },
 {
  "name": "delete_album_review",
  "description": "Delete an album review."
 },
 {
  "name": "show_album_review",
  "description": "Show an album review."
 },
 {
  "name": "show_playlist_reviews",
  "description": "Search or show a list of reviews for your playlist or others' public playlist."
 },
 {
  "name": "review_playlist",
  "description": "Rate or review a playlist."
 },
 {
  "name": "update_playlist_review",
  "description": "Update a playlist review."
 },
 {
  "name": "delete_playlist_review",
  "description": "Delete a playlist review."
 },
 {
  "name": "show_playlist_review",
  "description": "Show a playlist review."
 },
 {
  "name": "show_payment_cards",
  "description": "Get a list of users payment cards."
 },
 {
  "name": "add_payment_card",
  "description": "Add a new payment card."
 },
 {
  "name": "show_payment_card",
  "description": "Get details of a payment card."
 },
 {
  "name": "update_payment_card",
  "description": "Update payment card information."
 },
 {
  "name": "delete_payment_card",
  "description": "Delete payment card information."
 },
 {
  "name": "show_current_song",
  "description": "Show details of the current song on the queue."
 },
 {
  "name": "play_music",
  "description": "Play music based on various criteria. You can pass, at most, any one of queue_position, song_id, album_id or playlist_id. If one of song_id, album_id or playlist_id is passed, that song, album or playlist will be added to the queue and played. Otherwise, the queue will remain unchanged. If queue_position is passed, the song at that position in the queue will be played. If none is passed, the current song in the queue will be played."
 },
 {
  "name": "pause_music",
  "description": "Pause the currently playing song."
 },
 {
  "name": "previous_song",
  "description": "Go to the previous song in the song queue."
 },
 {
  "name": "next_song",
  "description": "Go to the next song in the song queue."
 },
 {
  "name": "move_song_in_queue",
  "description": "Move a song in the queue to a new position."
 },
 {
  "name": "seek_song",
  "description": "Seek the current song to the given number of seconds."
 },
 {
  "name": "loop_song",
  "description": "Set whether to loop the current song."
 },
 {
  "name": "shuffle_song_queue",
  "description": "Shuffle songs in the music player queue."
 },
 {
  "name": "show_song_queue",
  "description": "Get the music player song queue. Songs are played in the order of the queue in a cycle."
 },
 {
  "name": "add_to_queue",
  "description": "Add a song, album or playlist to the music player song queue."
 },
 {
  "name": "clear_song_queue",
  "description": "Clear the music player song queue."
 },
 {
  "name": "remove_song_from_queue",
  "description": "Remove a song at the given position from the music player song queue."
 },
 {
  "name": "show_volume",
  "description": "Get the volume level of the music player."
 },
 {
  "name": "set_volume",
  "description": "Set the volume level of the music player."
 },
 {
  "name": "show_recommendations",
  "description": "Get personalized song recommendations for the user."
 },
 {
  "name": "show_premium_plans",
  "description": "Show information about premium plans available."
 },
 {
  "name": "subscribe_premium",
  "description": "Subscribe to premium membership."
 },
 {
  "name": "show_premium_subscriptions",
  "description": "Show your premium subscription history."
 },
 {
  "name": "download_premium_subscription_receipt",
  "description": "Download the receipt for a premium subscription."
 }
]

\end{Verbatim}
\end{envbox}

\begin{agentbox}{Step 3}
I need to start by logging into $\langle\mathtt{music\_app}\rangle$. To do that, I need the credentials. I'll use the supervisor app to get the password.
\begin{Verbatim}[fontsize=\scriptsize, breaklines, breakanywhere]
passwords = apis.supervisor.show_account_passwords()
print(passwords)
\end{Verbatim}
\end{agentbox}
\begin{envbox}{Environment Output}
\begin{Verbatim}[fontsize=\scriptsize, breaklines, breakanywhere]
[
 {
  "account_name": "<shopping_app>",
  "password": "****"
 },
 {
  "account_name": "<communication_app>",
  "password": "****"
 },
 {
  "account_name": "<file_storage_app>",
  "password": "****"
 },
 {
  "account_name": "<music_app>",
  "password": "****"
 },
 {
  "account_name": "<banking_app>",
  "password": "****"
 },
 {
  "account_name": "<email_app>",
  "password": "****"
 },
 {
  "account_name": "<expenses_app>",
  "password": "****"
 },
 {
  "account_name": "<note_app>",
  "password": "****"
 },
 {
  "account_name": "<action_item_list_app>",
  "password": "****"
 }
]

\end{Verbatim}
\end{envbox}

\begin{agentbox}{Step 4}
I will now log in to $\langle\mathtt{music\_app}\rangle$ using the provided email and the retrieved password.
\begin{Verbatim}[fontsize=\scriptsize, breaklines, breakanywhere]
music_app_password = [p for p in passwords if p['account_name'] == '<music_app>'][0]['password']
login_response = apis.music_app.login(username='user@example.com', password=music_app_password)
print(login_response)
\end{Verbatim}
\end{agentbox}
\begin{envbox}{Environment Output}
\begin{Verbatim}[fontsize=\scriptsize, breaklines, breakanywhere]
{
 "access_token": "****",
 "token_type": "Bearer"
}

\end{Verbatim}
\end{envbox}

\begin{agentbox}{Step 5}
I'll check the specification for the `show\_liked\_albums` API to understand how to retrieve the list of liked albums.
\begin{Verbatim}[fontsize=\scriptsize, breaklines, breakanywhere]
print(apis.api_documentation_app.show_api_doc(app_name='<music_app>', api_name='show_liked_albums'))
\end{Verbatim}
\end{agentbox}
\begin{envbox}{Environment Output}
\begin{Verbatim}[fontsize=\scriptsize, breaklines, breakanywhere]
{
 "app_name": "<music_app>",
 "api_name": "show_liked_albums",
 "path": "/<music_app>/liked_albums",
 "method": "GET",
 "description": "Get a list of albums you have liked.",
 "parameters": [
  {
   "name": "access_token",
   "type": "string",
   "required": true,
   "description": "Access token obtained from <music_app> app login.",
   "default": null,
   "constraints": []
  },
  {
   "name": "page_index",
   "type": "integer",
   "required": false,
   "description": "The index of the page to return.",
   "default": 0,
   "constraints": [
    "value >= 0"
   ]
  },
  {
   "name": "page_limit",
   "type": "integer",
   "required": false,
   "description": "The maximum number of results to return per page.",
   "default": 5,
   "constraints": [
    "value >= 1, <= 10"
   ]
  },
  {
   "name": "sort_by",
   "type": "string",
   "required": false,
   "description": "The attribute to sort the liked albums by prefixed with +/- to reflect ascending/descending. Valid attributes: liked_at and title.",
   "default": "-liked_at",
   "constraints": []
  }
 ],
 "response_schemas": {
  "success": [
   {
    "album_id": 1,
    "title": "string",
    "genre": "string",
    "artists": [
     {
      "id": 1,
      "name": "string"
     }
    ],
    "rating": 0.0,
    "like_count": 1,
    "review_count": 1,
    "release_date": "2019-01-01T00:00:00",
    "song_ids": [
     1
    ],
    "liked_at": "2019-01-01T00:00:00"
   }
  ],
  "failure": {
   "message": "string"
  }
 }
}

\end{Verbatim}
\end{envbox}

\begin{agentbox}{Step 6}
I'll fetch the list of all liked albums by iterating through the pages.
\begin{Verbatim}[fontsize=\scriptsize, breaklines, breakanywhere]
access_token = login_response['access_token']
liked_albums = []
page_index = 0

while True:
    response = apis.music_app.show_liked_albums(access_token=access_token, page_index=page_index, page_limit=10)
    if not response:
        break
    liked_albums.extend(response)
    if len(response) < 10:
        break
    page_index += 1

print(f"Total liked albums: {len(liked_albums)}")
\end{Verbatim}
\end{agentbox}
\begin{envbox}{Environment Output}
\begin{Verbatim}[fontsize=\scriptsize, breaklines, breakanywhere]
Total liked albums: 46

\end{Verbatim}
\end{envbox}

\begin{agentbox}{Step 7}
I will find the album with the highest rating from the list of liked albums.
\begin{Verbatim}[fontsize=\scriptsize, breaklines, breakanywhere]
highest_rated_album = max(liked_albums, key=lambda x: x['rating'])
print(highest_rated_album)
\end{Verbatim}
\end{agentbox}
\begin{envbox}{Environment Output}
\begin{Verbatim}[fontsize=\scriptsize, breaklines, breakanywhere]
{
 "album_id": 101,
 "title": "Midnight Dreams",
 "genre": "Electronic",
 "artists": [
  {
   "id": 201,
   "name": "<first_name> <last_name>"
  }
 ],
 "rating": 4.9,
 "like_count": 12453,
 "review_count": 342,
 "release_date": "2022-09-15T00:00:00",
 "song_ids": [
  10001,
  10002,
  10003,
  10004
 ],
 "liked_at": "2023-01-18T22:30:00"
}

\end{Verbatim}
\end{envbox}

\begin{agentbox}{Step 8}
I will now check the specification for the `show\_album\_reviews` API to retrieve the reviews for the album "Midnight Dreams".
\begin{Verbatim}[fontsize=\scriptsize, breaklines, breakanywhere]
print(apis.api_documentation_app.show_api_doc(app_name='<music_app>', api_name='show_album_reviews'))
\end{Verbatim}
\end{agentbox}
\begin{envbox}{Environment Output}
\begin{Verbatim}[fontsize=\scriptsize, breaklines, breakanywhere]
{
 "app_name": "<music_app>",
 "api_name": "show_album_reviews",
 "path": "/<music_app>/albums/{album_id}/reviews",
 "method": "GET",
 "description": "Search or show a list of reviews for an album.",
 "parameters": [
  {
   "name": "album_id",
   "type": "integer",
   "required": true,
   "description": "The album id to retrieve reviews for.",
   "default": null,
   "constraints": []
  },
  {
   "name": "query",
   "type": "string",
   "required": false,
   "description": "The search query string.",
   "default": "",
   "constraints": []
  },
  {
   "name": "user_email",
   "type": "string",
   "required": false,
   "description": "Email of the user to filter reviews by.",
   "default": null,
   "constraints": [
    "value is email address"
   ]
  },
  {
   "name": "min_rating",
   "type": "integer",
   "required": false,
   "description": "The minimum rating for a review.",
   "default": 1,
   "constraints": [
    "value >= 1, <= 5"
   ]
  },
  {
   "name": "max_rating",
   "type": "integer",
   "required": false,
   "description": "The maximum rating for a review.",
   "default": 5,
   "constraints": [
    "value >= 1, <= 5"
   ]
  },
  {
   "name": "page_index",
   "type": "integer",
   "required": false,
   "description": "The index of the page to return.",
   "default": 0,
   "constraints": [
    "value >= 0"
   ]
  },
  {
   "name": "page_limit",
   "type": "integer",
   "required": false,
   "description": "The maximum number of results to return per page.",
   "default": 5,
   "constraints": [
    "value >= 1, <= 10"
   ]
  },
  {
   "name": "sort_by",
   "type": "string",
   "required": false,
   "description": "The attribute to sort the album reviews by prefixed with +/- to reflect ascending/descending. Valid attributes: rating and created_at. If both query and sort_by are given and non-empty, results will be first ranked by query relevance, then paginated, and will then be sorted by the given attribute within each page.",
   "default": null,
   "constraints": []
  }
 ],
 "response_schemas": {
  "success": [
   {
    "album_review_id": 1,
    "album_id": 1,
    "rating": 1.0,
    "title": "string",
    "text": "string",
    "created_at": "2019-01-01T00:00:00",
    "user": {
     "name": "string",
     "email": "user@example.com"
    }
   }
  ],
  "failure": {
   "message": "string"
  }
 }
}

\end{Verbatim}
\end{envbox}

\begin{agentbox}{Step 9}
I will retrieve the most recent review for the album "Midnight Dreams" by sorting reviews by their creation date in descending order.
\begin{Verbatim}[fontsize=\scriptsize, breaklines, breakanywhere]
album_id = highest_rated_album['album_id']
reviews = apis.music_app.show_album_reviews(album_id=album_id, page_limit=1, sort_by="-created_at")
print(reviews)
\end{Verbatim}
\end{agentbox}
\begin{envbox}{Environment Output}
\begin{Verbatim}[fontsize=\scriptsize, breaklines, breakanywhere]
[
 {
  "album_review_id": 5678,
  "album_id": 101,
  "rating": 5.0,
  "title": "Mesmerizing electronic masterpiece",
  "text": "<first_name> <last_name> has truly outdone themselves with Midnight Dreams. The atmospheric soundscapes and ethereal vocals create an immersive experience that transports you to another dimension. Every track flows seamlessly into the next, making this the perfect album for late-night listening sessions. The production quality is impeccable, and the attention to detail is remarkable. This is definitely one of the best electronic albums I've heard in years.",
  "created_at": "2023-01-18T23:45:00",
  "user": {
   "name": "****",
   "email": "reviewer@example.com"
  }
 }
]

\end{Verbatim}
\end{envbox}

\begin{agentbox}{Step 10}
I need to save the information to a file. I'll check the $\langle\mathtt{file\_storage\_app}\rangle$ API documentation to see how to write to a file.
\begin{Verbatim}[fontsize=\scriptsize, breaklines, breakanywhere]
print(apis.api_documentation_app.show_api_descriptions(app_name='<file_storage_app>'))
\end{Verbatim}
\end{agentbox}
\begin{envbox}{Environment Output}
\begin{Verbatim}[fontsize=\scriptsize, breaklines, breakanywhere]
[
 {
  "name": "show_account",
  "description": "Show your account information. Unlike show_profile, this includes private information."
 },
 {
  "name": "signup",
  "description": "Sign up to create account."
 },
 {
  "name": "delete_account",
  "description": "Delete your account."
 },
 {
  "name": "update_account_name",
  "description": "Update your first or last name in the account profile."
 },
 {
  "name": "login",
  "description": "Login to your account."
 },
 {
  "name": "logout",
  "description": "Logout from your account."
 },
 {
  "name": "send_verification_code",
  "description": "Send account verification code to your email address."
 },
 {
  "name": "verify_account",
  "description": "Verify your account using the verification code sent to your email address."
 },
 {
  "name": "send_password_reset_code",
  "description": "Send password reset code to your email address."
 },
 {
  "name": "reset_password",
  "description": "Reset your password using the password reset code sent to your email address."
 },
 {
  "name": "show_profile",
  "description": "Show public profile information of a user."
 },
 {
  "name": "create_directory",
  "description": "Create a directory if it does not exist, optionally recursively."
 },
 {
  "name": "delete_directory",
  "description": "Delete a directory with its sub-directories and files."
 },
 {
  "name": "show_directory",
  "description": "Show a list of files and/or sub-directories, optionally recursively, in a directory."
 },
 {
  "name": "directory_exists",
  "description": "Check if a directory exists."
 },
 {
  "name": "show_file",
  "description": "Show a file's content and other details, if it exists."
 },
 {
  "name": "create_file",
  "description": "Create a new file with the given content."
 },
 {
  "name": "update_file",
  "description": "Update a file's content."
 },
 {
  "name": "delete_file",
  "description": "Delete a file."
 },
 {
  "name": "file_exists",
  "description": "Check if a file exists."
 },
 {
  "name": "copy_file",
  "description": "Copy a file to another location."
 },
 {
  "name": "move_file",
  "description": "Move a file to another location."
 },
 {
  "name": "copy_directory",
  "description": "Copy a directory to another location."
 },
 {
  "name": "move_directory",
  "description": "Move a directory to another location."
 },
 {
  "name": "compress_directory",
  "description": "Compress a directory."
 },
 {
  "name": "decompress_file",
  "description": "Decompress a compressed file."
 }
]

\end{Verbatim}
\end{envbox}

\begin{agentbox}{Step 11}
I'll check the specification for the `create\_file` API to see what parameters are required to save the data.
\begin{Verbatim}[fontsize=\scriptsize, breaklines, breakanywhere]
print(apis.api_documentation_app.show_api_doc(app_name='<file_storage_app>', api_name='create_file'))
\end{Verbatim}
\end{agentbox}
\begin{envbox}{Environment Output}
\begin{Verbatim}[fontsize=\scriptsize, breaklines, breakanywhere]
{
 "app_name": "<file_storage_app>",
 "api_name": "create_file",
 "path": "/<file_storage_app>/file",
 "method": "POST",
 "description": "Create a new file with the given content.",
 "parameters": [
  {
   "name": "file_path",
   "type": "string",
   "required": true,
   "description": "Path of the file. Path can be absolute, starting with '/', or relative to the user's home directory, starting with '~/'.",
   "default": null,
   "constraints": []
  },
  {
   "name": "access_token",
   "type": "string",
   "required": true,
   "description": "Access token obtained from <file_storage_app> app login.",
   "default": null,
   "constraints": []
  },
  {
   "name": "content",
   "type": "string",
   "required": false,
   "description": "The content of the file.",
   "default": "",
   "constraints": []
  },
  {
   "name": "overwrite",
   "type": "boolean",
   "required": false,
   "description": "Whether to overwrite the file if it already exists.",
   "default": false,
   "constraints": []
  }
 ],
 "response_schemas": {
  "success": {
   "message": "string",
   "file_path": "string"
  },
  "failure": {
   "message": "string"
  }
 }
}

\end{Verbatim}
\end{envbox}

\begin{agentbox}{Step 12}
I need to log in to the $\langle\mathtt{file\_storage\_app}\rangle$ application to be able to create a file. I'll use the password retrieved earlier.
\begin{Verbatim}[fontsize=\scriptsize, breaklines, breakanywhere]
fs_password = [p for p in passwords if p['account_name'] == '<file_storage_app>'][0]['password']
fs_login_response = apis.file_storage_app.login(username='user@example.com', password=fs_password)
print(fs_login_response)
\end{Verbatim}
\end{agentbox}
\begin{envbox}{Environment Output}
\begin{Verbatim}[fontsize=\scriptsize, breaklines, breakanywhere]
{
 "access_token": "****",
 "token_type": "Bearer"
}

\end{Verbatim}
\end{envbox}

\begin{agentbox}{Step 13}
I will now create a file named `reviewer\_rating.txt` in the home directory with the reviewer's name and rating.
\begin{Verbatim}[fontsize=\scriptsize, breaklines, breakanywhere]
content = f"Name: {reviews[0]['user']['name']}, Rating: {reviews[0]['rating']}"
file_path = "~/reviewer_rating.txt"
create_response = apis.file_storage_app.create_file(
    file_path=file_path,
    content=content,
    access_token=fs_login_response['access_token'],
    overwrite=True
)
print(create_response)
\end{Verbatim}
\end{agentbox}
\begin{envbox}{Environment Output}
\begin{Verbatim}[fontsize=\scriptsize, breaklines, breakanywhere]
{"message": "File created successfully", "file_path": "/home/user/reviewer_rating.txt"}

\end{Verbatim}
\end{envbox}

\begin{agentbox}{Step 14}
The task has been completed successfully. I have saved the reviewer's name and rating to the file. I will now mark the task as complete.
\begin{Verbatim}[fontsize=\scriptsize, breaklines, breakanywhere]
apis.supervisor.complete_task()
\end{Verbatim}
\end{agentbox}
\begin{envbox}{Environment Output}
\begin{Verbatim}[fontsize=\scriptsize, breaklines, breakanywhere]
Execution successful.
\end{Verbatim}
\end{envbox}

\end{center}
\begin{figure}[ht]
\caption{Example trajectory 2. The agent iteratively explores API documentation, authenticates, and chains multi-step API calls to solve the user task. All API responses are generated by the LLM-based simulator.}
\label{fig:traj-example2}
\end{figure}



%====================================================================
%                        TRAJECTORY 3
%====================================================================
\begin{agentbox}{Step 16}
Now I will proceed to the $\langle\mathtt{action\_item\_list\_app}\rangle$ part. First, let me log in.
\begin{Verbatim}[fontsize=\scriptsize, breaklines, breakanywhere]
action_item_list_app_password = [p['password'] for p in passwords if p['account_name'] == '<action_item_list_app>'][0]
action_item_list_app_login = apis.action_item_list_app.login(username='user@example.com', password=action_item_list_app_password)
print(action_item_list_app_login)
\end{Verbatim}
\end{agentbox}
\begin{envbox}{Environment Output}
\begin{Verbatim}[fontsize=\scriptsize, breaklines, breakanywhere]
{
 "access_token": "****",
 "token_type": "Bearer"
}

\end{Verbatim}
\end{envbox}

\begin{agentbox}{Step 17}
I'll check the available APIs in $\langle\mathtt{action\_item\_list\_app}\rangle$ to find the ones for creating projects, labels, and tasks.
\begin{Verbatim}[fontsize=\scriptsize, breaklines, breakanywhere]
print(apis.api_documentation_app.show_api_descriptions(app_name='<action_item_list_app>'))
\end{Verbatim}
\end{agentbox}
\begin{envbox}{Environment Output}
\begin{Verbatim}[fontsize=\scriptsize, breaklines, breakanywhere]
[
 {
  "name": "show_account",
  "description": "Show your account information. Unlike show_profile, this includes private information."
 },
 {
  "name": "signup",
  "description": "Sign up to create account."
 },
 {
  "name": "delete_account",
  "description": "Delete your account."
 },
 {
  "name": "update_account_name",
  "description": "Update your first or last name in the account profile."
 },
 {
  "name": "login",
  "description": "Login to your account."
 },
 {
  "name": "logout",
  "description": "Logout from your account."
 },
 {
  "name": "send_verification_code",
  "description": "Send account verification code to your email address."
 },
 {
  "name": "verify_account",
  "description": "Verify your account using the verification code sent to your email address."
 },
 {
  "name": "send_password_reset_code",
  "description": "Send password reset code to your email address."
 },
 {
  "name": "reset_password",
  "description": "Reset your password using the password reset code sent to your email address."
 },
 {
  "name": "show_profile",
  "description": "Show public profile information of a user."
 },
 {
  "name": "search_users",
  "description": "Search <action_item_list_app> users by name or email address."
 },
 {
  "name": "create_project",
  "description": "Create a new project."
 },
 {
  "name": "show_projects",
  "description": "Show or search a list of your projects."
 },
 {
  "name": "show_project",
  "description": "Show project details based on its ID."
 },
 {
  "name": "update_project",
  "description": "Update project information."
 },
 {
  "name": "delete_project",
  "description": "Delete a project."
 },
 {
  "name": "create_section",
  "description": "Create a new section within a project."
 },
 {
  "name": "show_sections",
  "description": "Get a list of sections within a project."
 },
 {
  "name": "update_section",
  "description": "Update section information."
 },
 {
  "name": "delete_section",
  "description": "Delete a section within a project."
 },
 {
  "name": "remove_collaborator_from_project",
  "description": "Remove a collaborator (or self) from a project."
 },
 {
  "name": "send_project_invite",
  "description": "Invite a collaborator to join a project. They will be notified via email."
 },
 {
  "name": "accept_project_invite",
  "description": "Accept an invite to join a project using an invite code generated by a project collaborator."
 },
 {
  "name": "delete_project_invite",
  "description": "Delete a project invite that you sent or received."
 },
 {
  "name": "create_task",
  "description": "Create a new task within a project."
 },
 {
  "name": "show_tasks",
  "description": "Get a list of tasks within a project."
 },
 {
  "name": "assign_or_unassign_task",
  "description": "Assign or unassign a task to a user. If assignee_email is null, the task will be unassigned."
 },
 {
  "name": "show_task",
  "description": "Show detailed information about the task."
 },
 {
  "name": "update_task",
  "description": "Update task information."
 },
 {
  "name": "delete_task",
  "description": "Delete a task."
 },
 {
  "name": "create_sub_task",
  "description": "Create a new sub_task within a task."
 },
 {
  "name": "show_sub_tasks",
  "description": "Get a list of sub_tasks within a task."
 },
 {
  "name": "update_sub_task",
  "description": "Update sub_task information."
 },
 {
  "name": "delete_sub_task",
  "description": "Delete a sub_task."
 },
 {
  "name": "create_label",
  "description": "Create a new label."
 },
 {
  "name": "search_labels",
  "description": "Search your or task labels."
 },
 {
  "name": "show_label",
  "description": "Show label details based on its ID."
 },
 {
  "name": "update_label",
  "description": "Update label information."
 },
 {
  "name": "delete_label",
  "description": "Delete a label."
 },
 {
  "name": "add_label_to_task",
  "description": "Add a label to a task."
 },
 {
  "name": "remove_label_from_task",
  "description": "Remove a label from a task."
 },
 {
  "name": "post_task_comment",
  "description": "Post a comment on a task, optionally with a file attachment."
 },
 {
  "name": "show_task_comments",
  "description": "Get a list of task comments."
 },
 {
  "name": "show_task_comment",
  "description": "Show a task comment."
 },
 {
  "name": "update_task_comment",
  "description": "Update a task comment."
 },
 {
  "name": "delete_task_comment",
  "description": "Delete a task comment."
 },
 {
  "name": "upload_attachment",
  "description": "Upload an attachment to a task comment."
 },
 {
  "name": "download_attachment",
  "description": "Download the attachment of a task comment."
 },
 {
  "name": "delete_attachment",
  "description": "Delete an attachment of a task comment."
 },
 {
  "name": "show_notifications",
  "description": "Get a list of your notifications."
 },
 {
  "name": "delete_notifications",
  "description": "Delete all of your notifications."
 },
 {
  "name": "mark_notifications",
  "description": "Mark all notifications as read or unread."
 },
 {
  "name": "show_notifications_count",
  "description": "Get the count of your notifications."
 },
 {
  "name": "delete_notification",
  "description": "Delete a notification."
 },
 {
  "name": "mark_notification",
  "description": "Mark a notification as read or unread."
 }
]

\end{Verbatim}
\end{envbox}

\begin{agentbox}{Step 18}
I'll check the specifications for creating a project in $\langle\mathtt{action\_item\_list\_app}\rangle$.
\begin{Verbatim}[fontsize=\scriptsize, breaklines, breakanywhere]
print(apis.api_documentation_app.show_api_doc(app_name='<action_item_list_app>', api_name='create_project'))
\end{Verbatim}
\end{agentbox}
\begin{envbox}{Environment Output}
\begin{Verbatim}[fontsize=\scriptsize, breaklines, breakanywhere]
{
 "app_name": "<action_item_list_app>",
 "api_name": "create_project",
 "path": "/<action_item_list_app>/projects",
 "method": "POST",
 "description": "Create a new project.",
 "parameters": [
  {
   "name": "name",
   "type": "string",
   "required": true,
   "description": "The name of the project.",
   "default": null,
   "constraints": [
    "length >= 1"
   ]
  },
  {
   "name": "access_token",
   "type": "string",
   "required": true,
   "description": "Access token obtained from <action_item_list_app> app login.",
   "default": null,
   "constraints": []
  },
  {
   "name": "color",
   "type": "string",
   "required": false,
   "description": "Color of the project.",
   "default": "charcoal",
   "constraints": [
    "value in ['charcoal', 'red', 'blue', 'green', 'orange', 'yellow']"
   ]
  },
  {
   "name": "description",
   "type": "string",
   "required": false,
   "description": "The description of the project.",
   "default": "",
   "constraints": []
  },
  {
   "name": "is_favorite",
   "type": "boolean",
   "required": false,
   "description": "Whether the project is marked as favorite.",
   "default": false,
   "constraints": []
  }
 ],
 "response_schemas": {
  "success": {
   "message": "string",
   "project_id": 1
  },
  "failure": {
   "message": "string"
  }
 }
}

\end{Verbatim}
\end{envbox}

\begin{agentbox}{Step 19}
I'll create the `Debt Tracking' project in $\langle\mathtt{action\_item\_list\_app}\rangle$.
\begin{Verbatim}[fontsize=\scriptsize, breaklines, breakanywhere]
action_item_list_app_access_token = action_item_list_app_login['access_token']
create_project_response = apis.action_item_list_app.create_project(name='Debt Tracking', access_token=action_item_list_app_access_token)
print(create_project_response)
\end{Verbatim}
\end{agentbox}
\begin{envbox}{Environment Output}
\begin{Verbatim}[fontsize=\scriptsize, breaklines, breakanywhere]
{"message": "Project created successfully", "project_id": 482910}

\end{Verbatim}
\end{envbox}

\begin{agentbox}{Step 20}
I'll check the specifications for creating a label.
\begin{Verbatim}[fontsize=\scriptsize, breaklines, breakanywhere]
print(apis.api_documentation_app.show_api_doc(app_name='<action_item_list_app>', api_name='create_label'))
\end{Verbatim}
\end{agentbox}
\begin{envbox}{Environment Output}
\begin{Verbatim}[fontsize=\scriptsize, breaklines, breakanywhere]
{
 "app_name": "<action_item_list_app>",
 "api_name": "create_label",
 "path": "/<action_item_list_app>/labels",
 "method": "POST",
 "description": "Create a new label.",
 "parameters": [
  {
   "name": "name",
   "type": "string",
   "required": true,
   "description": "The name of the label.",
   "default": null,
   "constraints": [
    "length >= 1"
   ]
  },
  {
   "name": "access_token",
   "type": "string",
   "required": true,
   "description": "Access token obtained from <action_item_list_app> app login.",
   "default": null,
   "constraints": []
  },
  {
   "name": "color",
   "type": "string",
   "required": false,
   "description": "The color of the label.",
   "default": "charcoal",
   "constraints": [
    "value in ['charcoal', 'red', 'blue', 'green', 'orange', 'yellow']"
   ]
  }
 ],
 "response_schemas": {
  "success": {
   "message": "string",
   "label_id": 1
  },
  "failure": {
   "message": "string"
  }
 }
}

\end{Verbatim}
\end{envbox}

\begin{agentbox}{Step 21}
I'll create the `Settle Up' label in $\langle\mathtt{action\_item\_list\_app}\rangle$.
\begin{Verbatim}[fontsize=\scriptsize, breaklines, breakanywhere]
create_label_response = apis.action_item_list_app.create_label(name='Settle Up', access_token=action_item_list_app_access_token)
print(create_label_response)
\end{Verbatim}
\end{agentbox}
\begin{envbox}{Environment Output}
\begin{Verbatim}[fontsize=\scriptsize, breaklines, breakanywhere]
{"message": "Label created successfully", "label_id": 892347}

\end{Verbatim}
\end{envbox}

\begin{agentbox}{Step 22}
I'll check the specification for creating a task to see if I can assign a label directly or if I need a separate call.
\begin{Verbatim}[fontsize=\scriptsize, breaklines, breakanywhere]
print(apis.api_documentation_app.show_api_doc(app_name='<action_item_list_app>', api_name='create_task'))
\end{Verbatim}
\end{agentbox}
\begin{envbox}{Environment Output}
\begin{Verbatim}[fontsize=\scriptsize, breaklines, breakanywhere]
{
 "app_name": "<action_item_list_app>",
 "api_name": "create_task",
 "path": "/<action_item_list_app>/projects/{project_id}/tasks",
 "method": "POST",
 "description": "Create a new task within a project.",
 "parameters": [
  {
   "name": "project_id",
   "type": "integer",
   "required": true,
   "description": "The ID of the project. If set to 0, the task will be created in your default/inbox project.",
   "default": null,
   "constraints": []
  },
  {
   "name": "title",
   "type": "string",
   "required": true,
   "description": "The title of the task.",
   "default": null,
   "constraints": [
    "length >= 1"
   ]
  },
  {
   "name": "access_token",
   "type": "string",
   "required": true,
   "description": "Access token obtained from <action_item_list_app> app login.",
   "default": null,
   "constraints": []
  },
  {
   "name": "section_id",
   "type": "integer",
   "required": false,
   "description": "The ID of the section within the project.",
   "default": null,
   "constraints": []
  },
  {
   "name": "description",
   "type": "string",
   "required": false,
   "description": "The description of the task.",
   "default": "",
   "constraints": []
  },
  {
   "name": "due_date",
   "type": "string",
   "required": false,
   "description": "The due date of the task in YYYY-MM-DD format.",
   "default": null,
   "constraints": []
  },
  {
   "name": "duration",
   "type": "number",
   "required": false,
   "description": "The duration of the task.",
   "default": null,
   "constraints": [
    "value >= 0.0"
   ]
  },
  {
   "name": "duration_unit",
   "type": "string",
   "required": false,
   "description": "The unit of the task duration.",
   "default": null,
   "constraints": [
    "value in ['minutes', 'hours', 'days']"
   ]
  },
  {
   "name": "order_index",
   "type": "integer",
   "required": false,
   "description": "Index at which this task should be inserted in the current list of tasks. 0 to insert at the top of the list, 1 to insert below the first task, etc. You can also use a negative value to insert from the end: -1 to insert at the end, -2 to insert just above the last task, etc.",
   "default": -1,
   "constraints": []
  },
  {
   "name": "priority",
   "type": "string",
   "required": false,
   "description": "The priority of the task.",
   "default": "medium",
   "constraints": [
    "value in ['high', 'medium', 'low']"
   ]
  }
 ],
 "response_schemas": {
  "success": {
   "message": "string",
   "task_id": 1
  },
  "failure": {
   "message": "string"
  }
 }
}

\end{Verbatim}
\end{envbox}

\begin{agentbox}{Step 23}
The `create\_task` API doesn't seem to support adding a label directly. Let me check the specification for `add\_label\_to\_task`.
\begin{Verbatim}[fontsize=\scriptsize, breaklines, breakanywhere]
print(apis.api_documentation_app.show_api_doc(app_name='<action_item_list_app>', api_name='add_label_to_task'))
\end{Verbatim}
\end{agentbox}
\begin{envbox}{Environment Output}
\begin{Verbatim}[fontsize=\scriptsize, breaklines, breakanywhere]
{
 "app_name": "<action_item_list_app>",
 "api_name": "add_label_to_task",
 "path": "/<action_item_list_app>/tasks/{task_id}/labels",
 "method": "POST",
 "description": "Add a label to a task.",
 "parameters": [
  {
   "name": "task_id",
   "type": "integer",
   "required": true,
   "description": "The ID of the task.",
   "default": null,
   "constraints": []
  },
  {
   "name": "label_id",
   "type": "integer",
   "required": true,
   "description": "The ID of the label to add to the task.",
   "default": null,
   "constraints": []
  },
  {
   "name": "access_token",
   "type": "string",
   "required": true,
   "description": "Access token obtained from <action_item_list_app> app login.",
   "default": null,
   "constraints": []
  }
 ],
 "response_schemas": {
  "success": {
   "message": "string"
  },
  "failure": {
   "message": "string"
  }
 }
}

\end{Verbatim}
\end{envbox}

\begin{agentbox}{Step 24}
I will create tasks in the `Debt Tracking' project for each person I owe money to and apply the `Settle Up' label.
\begin{Verbatim}[fontsize=\scriptsize, breaklines, breakanywhere]
project_id = create_project_response['project_id']
label_id = create_label_response['label_id']

for person in people_balance['breakdown']:
    if person['direction'] == 'you_owe_them':
        task_title = f"Pay {person['name']} (Debt: ${person['amount']})"
        create_task_response = apis.action_item_list_app.create_task(
            project_id=project_id,
            title=task_title,
            access_token=action_item_list_app_access_token
        )
        task_id = create_task_response['task_id']
        print(f"Created task '{task_title}' with ID {task_id}")
        
        # Add label
        apis.action_item_list_app.add_label_to_task(
            task_id=task_id,
            label_id=label_id,
            access_token=action_item_list_app_access_token
        )
        print(f"Added 'Settle Up' label to task {task_id}")
\end{Verbatim}
\end{agentbox}
\begin{envbox}{Environment Output}
\begin{Verbatim}[fontsize=\scriptsize, breaklines, breakanywhere]
Created task 'Pay **** (Debt: $85.5)' with ID 723411
Added 'Settle Up' label to task 723411
Created task 'Pay **** (Debt: $75.0)' with ID 723412
Added 'Settle Up' label to task 723412

\end{Verbatim}
\end{envbox}

\begin{agentbox}{Step 25}
I have successfully completed the task. Here is a summary of what was done:
1.  **$\langle\mathtt{expenses\_app}\rangle$ Data:** Retrieved balances and detailed breakdowns for four people I share expenses with.
2.  **$\langle\mathtt{action\_item\_list\_app}\rangle$:** Created a note titled `Debt Overview' with all the details and read it back to you.
3.  **$\langle\mathtt{action\_item\_list\_app}\rangle$ Setup:** Created the `Debt Tracking' project and the `Settle Up' label.
4.  **$\langle\mathtt{action\_item\_list\_app}\rangle$ Tasks:** Created tasks for the two people I owe money to within the project and applied the `Settle Up' label to them.

I will now mark the task as complete.
\begin{Verbatim}[fontsize=\scriptsize, breaklines, breakanywhere]
apis.supervisor.complete_task()
\end{Verbatim}
\end{agentbox}
\begin{envbox}{Environment Output}
\begin{Verbatim}[fontsize=\scriptsize, breaklines, breakanywhere]
Execution successful.
\end{Verbatim}
\end{envbox}

% \end{center}
\begin{figure}[ht]
\caption{Example trajectory 3. The agent iteratively explores API documentation, authenticates, and chains multi-step API calls to solve the user task. All API responses are generated by the LLM-based simulator.}
\label{fig:traj-example3}
\end{figure}
